\section{Lohengrin}

\begin{quotex}
This is a chapter from \textit{The Middle Ages} by \textbf{Karl Heyer}.
\end{quotex}


In the knowledge of the Mysteries it is known that the great leaders of mankind, the initiates, are those who help humanity take a new step in development. Lohengrin is the messenger of the Holy Grail. He appears before medieval consciousness as the great leader and initiate who takes mankind a step further during the middle part of the Middle Ages. He brought the culture of the towns, he inspired the citizenry as they came into being. This is the Lohengrin individuality: and Elsa of Brabant is nothing but a symbol of the medieval folk soul who is to take a further step of development under the influence of Lohengrin. The legend of this progress in human history is beautifully and movingly portrayed.

By the question she asks from curiosity, Elsa brings about the disappearance of the initiate. This expresses the increasing earthly darkening of consciousness of the people who find their path of development in the towns. Materialism sets in, and only centuries later will this darkened consciousness, that has, however, set human beings free, once more receive a call from the spiritual world, a call for further great changes of consciousness coming from the continuing, though transformed stream of the Holy Grail.

The fertile seeds of the modern economic realm, as it later grew on capitalist, bourgeois foundations during the centuries of the fifth post-Atlantean era, were certainly sown in the Middle Ages. The same cannot be said, however, of the modern state. Its roots lie the territories of the princes who gradually began to organise their states with the help of legally trained civil servants. This tendency was based on the countryside, not on the towns which were often engaged in direct dispute with the princes. With regard to the time beginning around the fifteenth century, Rudolf Steiner pointed to the following:

\begin{quotex}
People in the towns were proud of their individuality and freedom, as may be seen in the portraits that have come down to us from those days... But the village communities remained outside. The power of the territorial princes made itself felt. People in the villages who gradually found themselves in opposition to the towns found their leaders among the ones who took their part, or said they would take their part, against the towns.
\end{quotex}


Under the influence of pressure from the surrounding countryside the towns became a part of wider administrative structures that embodied the principles of Roman law. The modern state came into being. It was formed working inwards towards the towns from the surrounding countryside. What conquered the towns from the surrounding countryside was a juridical Latin or Roman element that was beginning to emerge and had indeed grown so powerful that there was no longer any chance for another stream that was struggling to come to the surface among the rural populations. In England and Bohemia, for example this was the stream of Wycliffe and Huss, and the Bohemian Brethren. None of this managed to take hold. The only thing that really succeeded in gaining the upper hand was the Roman administrative element.


\flrightit{Posted on 2022-09-07 by Cologero}

